\documentclass[10pt,a4paper]{amsart}
\usepackage[margin=19mm]{geometry}
\usepackage{amsmath,amssymb,mathtools}
\usepackage[scr=rsfs,cal=euler]{mathalfa}
\usepackage{xfrac}
\usepackage[unicode,hidelinks,bookmarks=false]{hyperref}
\newcommand{\ie}{i.e.\ }
\newcommand{\cf}{cf.\ }
\newcommand{\resp}{respectively\ }
\newcommand{\Subsection}{\subsection}
\providecommand{\nobreakdash}{\nobreak\hbox{-}}
\setlength{\emergencystretch}{2em}
\multlinegap0pt
\usepackage{bm}
\usepackage{bbm}
\usepackage{dsfont}
\usepackage{xspace}
\usepackage{cancel}
\usepackage{mathtools}
\usepackage{amsthm}
\makeatletter
\@ifundefined{theorem}{\newtheorem{theorem}{Theorem}}{}
\@ifundefined{lemma}{\newtheorem{lemma}[theorem]{Lemma}}{}
\@ifundefined{proposition}{\newtheorem{proposition}[theorem]{Proposition}}{}
\makeatother
\def \O{\Omega}
\title[Variational method for fractional Hamiltonian system in boun]{Variational method for fractional Hamiltonian system in bounded domain}
\author{Weimin Zhang}
\date{Source: arXiv:2404.00687v2 (CC BY 4.0)}
\begin{document}
\maketitle

\noindent\emph{Source and licence.} Variational method for fractional Hamiltonian system in bounded domain. Version arXiv:2404.00687v2 (CC BY 4.0), distributed under the Creative Commons Attribution 4.0 International licence, as recorded in the arXiv licence metadata for this exact version and shown on the abstract page https://arxiv.org/abs/2404.00687v2. Full rights notice: licenses/arxiv-2404.00687-CC-BY-4.0.txt.\par\smallskip

\noindent\textit{Editorial scope.} Counted unit: the Proposition whose source label is \texttt{2302102144} in the article (source0.tex lines 756--772, source label \texttt{2302102144}), delivered together with the article's own complete original proof of it (source0.tex lines 774--810, one \texttt{proof} environment of 37 lines). No mathematics is written, repaired or re-derived: statement and proof are the article's own text line for line. Required context retained uncounted: 2306011244 (lines 352--355); 2402061816 (lines 366--372); 2304232341 (lines 707--712); 2211162013 (lines 735--751). Every definition, display and label of those lines is retained unchanged. Omitted from this package and not redistributed: 1--351, 356--365, 373--706, 713--734, 752--755, 773--773, 811--1786. Nothing inside a retained excerpt is omitted or altered: every retained body is delivered exactly as the article's lines are written, so no omission receipt of any kind accompanies this package. Citations: the retained excerpts carry no citation command at all, so no bibliography entry of the article is transcribed and no reference is redistributed. Reference adaptation: none was needed and none was made; every reference inside the delivered unit resolves inside it, so no unresolved reference or citation is produced. Printed slips kept literal: no printed word, symbol, hypothesis or number of the counted statement or of its proof was altered. Layout and class: the article's own class is \texttt{article} with packages including bbm, amsfonts, amsmath, amssymb, fancyhdr, setspace, latexsym, mathrsfs, amsthm, cite, enumitem, graphicx, hyperref, cleveref; the wrapper is the lane's pinned \texttt{amsart} editorial wrapper at 10pt on A4, which restates the theorem-like environments the retained text needs and adds the article's own macro definitions that the retained text needs (\texttt{\textbackslash O}), with extra wrapper packages (bm, bbm, dsfont, xspace, cancel, mathtools). The delivered numbering is the unit's own. Assets: no retained span contains a figure environment, a graphics-inclusion command, a PDF-page inclusion, a table or a TikZ picture; no image file is copied into this package, the package declares no asset dependency and no image file is redistributed with it. Licence: version arXiv:2404.00687v2 of this article is distributed under the Creative Commons Attribution 4.0 International licence, as recorded in the arXiv licence metadata for this exact version and shown on the abstract page https://arxiv.org/abs/2404.00687v2. A full rights notice is delivered at licenses/arxiv-2404.00687-CC-BY-4.0.txt. Characters: every retained excerpt is pure ASCII, so the delivered engine, XeLaTeX with its default Latin Modern text font, reports no missing character.
\begin{equation}\label{2306011244}
(-\Delta)^{s} u\;=f \;\;\text{in}~\Omega,\quad
u =0~\;\;\text{in} ~ \mathbb{R}^N\setminus\Omega.
\end{equation}

\begin{lemma}\label{2402061816}
Let $\Omega$ be a bounded $C^{1,1}$ domain and $f\in L^{\infty}(\Omega)$. There exists a unique $u\in X_0^s(\Omega)$ solving \eqref{2306011244}.  Moreover $u\in C^s(\overline{\Omega})\cap C_{\delta}^{\alpha}(\overline{\Omega})$, and
\[
\|u\|_{C^s(\overline{\Omega})}+ \left\|\frac{u}{\delta^s}\right\|_{C^\alpha(\overline{\Omega})}\le C\|f\|_{\infty}
\]
for some $0<\alpha<\min\{s, 1-s\}$. The constants $\alpha$, $C$ depend only on $\Omega$ and $s$.
\end{lemma}

\begin{proposition}\label{2304232341}
Given any $f\in L^1(\Omega; \delta^s dx) $, there exists a unique $L^1$-weak solution $u$ to \eqref{2306011244}. Moreover, there exists $C=C(s, \Omega)>0$ such that
\begin{equation}\label{2304142025}
|u|_1\le C\|f\|_{L^1(\Omega; \delta^s dx) }.
\end{equation}
\end{proposition}

\begin{proposition}\label{2211162013}
Assume that $s\in (0, 1)$, $N>2s$, $\Omega$ is a smooth bounded domain. For any $f\in L^r(\Omega)$ with $r\ge 1$, the unique $L^1$-weak solution $u$ to \eqref{2306011244} satisfies
\begin{itemize}
\item[\rm (i)] If $1 \le r < \frac{N}{2s}$, $u\in L^{\gamma}(\Omega)$ and $|u|_{\gamma}\le C(\Omega, r, s)|f|_r$,
where
\[
\gamma={\frac{Nr}{N-2rs}}\; {\rm if }\; r>1, \quad 1\le \gamma<\frac{N}{N-2s}\;{\rm if }\; r=1;
\]
\item[\rm (ii)] If $r \ge \frac{2N}{N+2s}$, we have $u\in X_0^s(\Omega)$ and $\|u\|_{X_0^s(\Omega)}\le C(\Omega, r, s)|f|_r$;
\item[\rm (iii)] If $r>\frac{N}{2s}$, there holds $u\in L^{\infty}(\Omega)$ and $|u|_{\infty}\le C(\Omega, r, s)|f|_{r}$.
\item[\rm (iv)] If $r=\frac{N}{2s}$, there is a constant $\alpha(\O, f, s) >0$ such that
\[
\int_{\Omega} e^{\alpha |u|} dx<\infty.
\]
In particular, $u\in L^\gamma(\Omega)$ for all $1\le \gamma<\infty$.
\end{itemize}
\end{proposition}

\begin{proposition}\label{2302102144}
Assume that $s\in (0,1)$, $N>2s$, $r\ge 1$, $\Omega$ is a smooth bounded domain. Then $\mathcal{W}^{2s, r}(\Omega)$ is a Banach space and $\mathcal{T}_s(\Omega)$ is dense in $\mathcal{W}^{2s, r}(\Omega)$. On the other hand, we have continuous embedding $\mathcal{W}^{2s, r}(\Omega)\subset L^{\gamma}(\Omega)$ for
\begin{equation}\label{2302110107}
\begin{cases}
\begin{aligned}
&1\le \gamma<\frac{N}{N-2s}, \;\; &&\text{if}~ r=1;\\
&1\le \gamma\le\frac{Nr}{N-2sr}, \;\; &&\text{if}~1 < r < \frac{N}{2s};\\
&1\le \gamma<\infty,  && \text{if}~2sr = N;\\
& \gamma = \infty,  && \text{if}~ 2sr > N.\\
\end{aligned}
\end{cases}
\end{equation}
Moreover, the above embeddings are compact provided
\begin{equation}\label{2302110109}
1\le \gamma<\frac{Nr}{N-2sr}\;\;  \text{if}~N>2sr;\quad \mbox{or} \;\; 1\le \gamma<\infty, \;\; \text{if}~N\le 2sr.
\end{equation}
\end{proposition}

\begin{proof}
The density of $\mathcal{T}_s(\Omega)$ in $\mathcal{W}^{2s, r}(\Omega)$ is obvious by definition and Lemma \ref{2402061816}. 

The completeness of $\mathcal{W}^{2s, r}(\Omega)$ is a simple fact. Let $\{u_j\}\subset \mathcal{W}^{2s, r}(\Omega)$ be a Cauchy sequence, by definition, $\{(-\Delta)^s u_j\}$ is a Cauchy sequence in $L^r({\Omega})$, hence converges to $v\in L^r({\Omega})$. Using Proposition \ref{2304232341}, there is a unique $L^1$-weak solution $u$ satisfying \eqref{2306011244} with $f = v$. Clearly $u \in \mathcal{W}^{2s, r}(\Omega)$, and $u_j$ tends to $u$ in $\mathcal{W}^{2s, r}(\Omega)$.

On the other hand, in virtue of Proposition \ref{2211162013}, the embedding of $\mathcal{W}^{2s, r}(\Omega)\subset L^\gamma(\Omega)$ is continuous if \eqref{2302110107} holds true, we shall only prove that the embedding is compact provided \eqref{2302110109}.

Without loss of generality, we only consider the case $N>2sr$, $1 < r$, $1 \le\gamma<\frac{Nr}{N-2sr}$. If $r\ge \frac{2N}{N+2s}$, then $\mathcal{W}^{2s, r}(\Omega)\subset X_0^s(\Omega)$ is continuous by Proposition \ref{2211162013}. By H\"older inequality, fix any $1 < \gamma<\frac{Nr}{N-2sr}$,
\begin{equation}\label{2302111630}
|u|_{{\gamma}}\le |u|^{1-\theta}_{{1}}|u|^{\theta}_{{\frac{Nr}{N-2sr}}} \quad \mbox{with } \theta=\frac{Nr-\gamma (N-2sr)}{Nr-(N-2sr)}.
\end{equation}
Since $\mathcal{W}^{2s, r}(\Omega)$ is continuously embedded in $L^{\frac{Nr}{N-2sr}}(\Omega)$ and compactly embedded  in $L^1(\Omega)$, $\mathcal{W}^{2s, r}(\Omega)$ is compactly embedded in $L^{\gamma}(\Omega)$.

Suppose now $1<r< \frac{2N}{N+2s}$, we first claim that the embedding $\mathcal{W}^{2s, r}(\Omega)\subset L^{1}(\Omega)$ is compact. Indeed, let $\{u_j\}$ be a bounded sequence in $\mathcal{W}^{2s, r}(\Omega)$, then $\{u_j\}$ is bounded in $L^{\frac{Nr}{N-2sr}}(\Omega)$ and the sequence $f_j = (-\Delta)^{s} u_j$ is bounded in $L^r(\Omega)$. Let $A > 0$ to be chosen later, we denote $g_j = f_j\chi_{|f_j(x)|\le A}$ and $h_j = f_j\chi_{|f_j(x)|> A}$. By H\"older and Chebychev inequalities,
\begin{align*}
\|h_j\|_{L^1(\Omega; \delta^s dx)} \leq C|h_j|_1\leq C|f_j|_r\Big(\int_{|f_j(x)|> A} dx\Big)^{1-\frac{1}{r}} \le C'A^{\frac{1}{r}-1}.
\end{align*}
Let $\{v_j\}$ and $\{w_j\}$ be respectively $L^1$-weak solutions of
\begin{equation}\label{2306010025}
\left\{
\begin{split}
(-\Delta)^{s} v_j &=g_j &&\text{in}~\Omega,\\
v_j&=0 &&\text{in} ~ \mathbb{R}^N\setminus\Omega,\\
\end{split}
\right. \quad \left\{
\begin{split}
(-\Delta)^{s} w_j&=h_j &&\text{in}~\Omega,\\
w_j&=0 &&\text{in} ~ \mathbb{R}^N\setminus\Omega.\\
\end{split}
\right.
\end{equation}
Clearly $u_j=v_j+w_j$. Given any $\varepsilon>0$, we fix $A$ large enough such that 
\[
|w_j|_1\le C\|h_j\|_{L^1(\Omega; \delta^s dx)} \le C'A^{\frac{1}{r}-1} \le \varepsilon.
\]
Moreover, $\{v_j\}$ is bounded in $X_0^s(\Omega)$ hence relatively compact in $L^1(\Omega)$. Therefore $\{u_j\}$ is precompact, or equivalently relatively compact in $L^1(\Omega)$, which means $\mathcal{W}^{2s, r}(\Omega)$ is compactly embedded in $L^{1}(\Omega)$. Applying again the interpolation inequality \eqref{2302111630}, $\{u_j\}$ is compact in $L^\gamma(\Omega)$ for $1\le\gamma<\frac{Nr}{N-2sr}$. So we are done.
\end{proof}

\end{document}
